\section{Utility gives no warning}
\label{sec:utility}

\paragraph{The held-out-loss optimum is a knife edge, not a margin.}
Early stopping is the one defence a practitioner already applies. At its own best held-out loss
every released model except the \kmer{3} reproduces \NUM{0\%} of the planted records, so stopping
exactly there does avoid verbatim extraction --- and that is the whole of the good news. The optimum
is almost flat in the loss and almost vertical in the leakage: one snapshot later costs HyenaDNA
\NUM{0.004}~\bits{} and takes it from \NUM{0\%} to \NUM{69\%} extraction, with the corresponding
gaps \NUM{0.032}~\bits{} to \NUM{100\%} for DNABERT-2, \NUM{0.039} to \NUM{51\%} for the \kmer{4},
\NUM{0.046} to \NUM{84\%} for the \kmer{5}, and \NUM{0.111} to \NUM{95\%} for the \kmer{6}
(Table~\ref{tab:stopping}). A validation loss within five hundredths of a bit of optimal is what
anyone would call converged, and inside that band sits every outcome from no leakage to complete
reproduction. The optimum is not clean on the other rulers either: recognition there is already
\NUM{4\%} for the \kmer{4}, \NUM{28\%} for the \kmer{5}, \NUM{52\%} for the \kmer{3}, \NUM{80\%} for
DNABERT-2 and \NUM{88\%} for the \kmer{6}. \textbf{Stopping at the optimum suppresses what the model
writes, not what it knows.}

\begin{table}[t]
\centering
\caption{What a practitioner's stopping rule selects. Next-token fine-tuning on the GUE promoter
set, three seeds; verbatim extraction of the $16\times$ canary at three candidate stopping points,
with the held-out loss or mean downstream MCC that would justify each choice.}
\label{tab:stopping}
\begin{tabular}{lccccc}
\toprule
& \multicolumn{2}{c}{best held-out loss} & \multicolumn{2}{c}{best downstream MCC} & last round \\
\cmidrule(lr){2-3}\cmidrule(lr){4-5}\cmidrule(lr){6-6}
model & round (loss) & extracted & round (MCC) & extracted & extracted \\
\midrule
DNABERT \kmer{3} & 30 (1.916) & 6\%  & 30 (0.655) & 6\%   & 37\% \\
DNABERT \kmer{4} & 15 (1.920) & 0\%  & 50 (0.663) & 51\%  & 51\% \\
DNABERT \kmer{5} & 15 (1.928) & 0\%  & 50 (0.654) & 84\%  & 84\% \\
DNABERT \kmer{6} & 15 (1.916) & 0\%  & 30 (0.663) & 44\%  & 95\% \\
DNABERT-2 (BPE)  &  3 (1.941) & 0\%  & 30 (0.642) & 100\% & 100\% \\
HyenaDNA medium  &  3 (1.913) & 0\%  & 50 (0.626) & 81\%  & 81\% \\
\bottomrule
\end{tabular}
\end{table}

\paragraph{Six downstream tasks do not warn either.}
At every snapshot we attach a classification head and fine-tune it on a GUE task
\cite{zhou2024dnabert2}: promoter detection, core promoters, transcription-factor binding and splice
sites in human, and the H3 and H3K4me3 histone marks in yeast. Over 50 rounds, while the planted
record goes from unmemorised to \NUM{37}--\NUM{100\%} verbatim reproduction, four of the six tasks
move by at most \NUM{0.037} MCC in either direction, which is within what we measure between seeds.
The two that move --- splice sites (\NUM{+0.02} to \NUM{+0.15}) and H3K4me3 (\NUM{+0.03} to
\NUM{+0.11}) --- move \emph{upward} as leakage grows, so a practitioner watching them is told to keep
training; transcription-factor binding drifts down for three models by at most \NUM{0.036}. None of
those three readings distinguishes a checkpoint at \NUM{0\%} extraction from the same model at
\NUM{95\%} (Appendix~A6).

\paragraph{Selecting by downstream score picks the leakiest checkpoint.}
For five of six models the best downstream mean is at round~30 or~50, where the leakage is:
\NUM{44}--\NUM{100\%} extraction, against \NUM{0}--\NUM{6\%} at the loss optimum of the same model
(Table~\ref{tab:stopping}). Only the \kmer{3} selects the same round under both rules. The frozen
linear probe behaves differently, peaking early and tracking the held-out loss, but it moves by at
most \NUM{0.03} MCC in total --- too little to act on. A practitioner selecting on benchmark
performance therefore lands close to the worst checkpoint for privacy, and one selecting on
validation loss is one snapshot away from it.

\paragraph{Probe validity.}
Our fine-tuning is a scaled-down version of the published recipe, so the absolute scores sit below
the published ones and only the movement within a column is ours to interpret. Five tasks land in
the published range; yeast H3K4me3 reaches only \NUM{0.22}--\NUM{0.30} and we read that row as a
weak probe rather than a reproduction. Four of the \NUM{252} (model, round, task) fine-tuning runs
collapsed to a degenerate head and are excluded and listed in Appendix~A6; that is the only
exclusion in this paper.
